\section{Related Work}
\label{sec:related}
\para{Self-reflection and iterative revision.} Methods such as Self-Refine and Reflexion show that language-only critique loops can improve outputs in many tasks \citep{madaan2023selfrefine,shinn2023reflexion}. These papers motivate reflection as a practical control strategy, but they focus on outcome metrics and do not directly test internal representation geometry.

\para{When intrinsic correction fails.} Several studies report that models often fail to detect their own reasoning errors, and that performance can degrade after self-correction unless stronger signals are provided \citep{huang2023llmscannot,chen2024errorlocation}. Tool-grounded critique (for example, CRITIC) can improve reliability by adding external verification \citep{gou2024critic}. Recent counter-evidence also argues that intrinsic correction can succeed under carefully controlled prompting and decoding conditions \citep{liu2024intrinsic}. Our work does not attempt to resolve all setting-dependent disagreements; instead, we test whether representation drift can distinguish critique-conditioned revision from generic second-pass generation.

\para{Representation geometry and activation-level control.} Representation engineering and activation steering show that behavior is often encoded in identifiable directions or subspaces \citep{zou2023representation,turner2023steering,arditi2024refusal}. Residual-stream analyses under knowledge conflict further demonstrate that intermediate states can predict response behavior \citep{zhao2024residual}. Building on this line, we focus on residual drift under reflection and explicitly compare it against a re-decode control to test whether self-critique creates a distinct geometric signature.

\para{Positioning.} Unlike prior reflection papers that stop at output gains, we jointly evaluate behavior, per-layer drift, and correction separability. Unlike prior representation papers that do not include reflection loops, we evaluate how critique-conditioned trajectories move relative to baseline and relative to a non-critique second attempt.
